% ============================================================================
% PREÁMBULO - Configuraciones y Paquetes
% ============================================================================

% Clase de documento
\documentclass[12pt,oneside]{book}

% -------------------------------------------------------------------------
% Paquetes fundamentales
% -------------------------------------------------------------------------

% Idioma español - Traducir "Cuadro" como "Tabla"
\usepackage[spanish]{babel}
\addto\captionsspanish{\def\tablename{Tabla}}

% Codificación UTF-8
\usepackage[utf8]{inputenc}

% Paquete para ecuaciones matemáticas
\usepackage{amsmath}
\usepackage{amssymb}

% Times New Roman (simula la fuente)
\usepackage{mathptmx}

% Interlineado 1.5
\usepackage{setspace}
\setstretch{1.5}

% Márgenes: Izquierdo 3.0 cm, Superior/Inferior/Derecho 2.5 cm
\usepackage{geometry}
\geometry{
    left=3.0cm,
    right=2.5cm,
    top=0.5cm,
    bottom=2.5cm
}

% -------------------------------------------------------------------------
% Paquetes adicionales
% -------------------------------------------------------------------------

% Control de líneas huérfanas/viudas
\widowpenalty=10000
\clubpenalty=10000

% Saltos de página controlados
\usepackage{emptypage}

% Formato de encabezados y pies de página
\usepackage{fancyhdr}
\fancyhf{}
\fancyhead[C]{}
\fancyfoot[C]{\thepage}
\pagestyle{fancy}

% Sangría francesa (primera línea)
\usepackage{indentfirst}

% URLs en bibliography
\usepackage{url}

% Hipervínculos en el PDF
\usepackage[colorlinks=true,linkcolor=black,anchorcolor=black,citecolor=black,filecolor=black,urlcolor=black]{hyperref}

% -------------------------------------------------------------------------
% Bibliografía - APA 7ma Edición (definida, uso manual)
% -------------------------------------------------------------------------
% Para habilitar citas automáticas: usar \cite{clave}
% \usepackage[backend=biber,style=apa,sorting=nyt]{biblatex}
% \addbibresource{referencias/bibliografia.bib}

% -------------------------------------------------------------------------
% Configuración de tablas y figuras
% -------------------------------------------------------------------------

\usepackage{graphicx}
\graphicspath{{figuras/}}

\usepackage{float}
\usepackage{array}
\usepackage{tabularx}
\usepackage{booktabs}
\usepackage[labelfont=bf]{caption}

% -------------------------------------------------------------------------
% Configuración de numeración de capítulos
% -------------------------------------------------------------------------

\usepackage{titlesec}

% Capítulo: "CAPÍTULO I"
\titleformat{\chapter}[display]
    {\bfseries\Large\centering}
    {\MakeUppercase{CAPÍTULO \Roman{chapter}}}
    {20pt}
    {\MakeUppercase}
    []

% Sección 1.1, 1.2 - justificada a la izquierda
\titleformat{\section}
    {\bfseries\large\raggedright}
    {\thesection}
    {0.5em}
    {}

% Subsección 1.1.1 - justificada a la izquierda
\titleformat{\subsection}
    {\bfseries\normalsize\raggedright}
    {\thesubsection}
    {0.5em}
    {}

% Paquete para bloques de código
\usepackage{listings}
\usepackage{xcolor}

\definecolor{codegreen}{rgb}{0,0.6,0}
\definecolor{codegray}{rgb}{0.5,0.5,0.5}
\definecolor{codepurple}{rgb}{0.58,0,0.82}
\definecolor{backcolour}{rgb}{0.96,0.96,0.98}

\lstdefinestyle{mystyle}{
    backgroundcolor=\color{backcolour},   
    commentstyle=\color{codegreen},
    keywordstyle=\color{magenta},
    numberstyle=\tiny\color{codegray},
    stringstyle=\color{codepurple},
    basicstyle=\ttfamily\footnotesize,
    breakatwhitespace=false,         
    breaklines=true,                 
    captionpos=b,                    
    keepspaces=true,                 
    numbers=left,                    
    numbersep=5pt,                  
    showspaces=false,                
    showstringspaces=false,
    showtabs=false,                  
    tabsize=2
}
\lstset{style=mystyle}